\chapter{Revision de antecedentes}

[Placeholder]

Revisión de antecedentes (ya sea productos, procesos, publicaciones, etc., a nivel académico o comercial) en el tema del trabajo. Puede incluir además (si es necesario) una breve introducción a los conceptos necesarios para entender el trabajo.

[Que poner]

En este capítulo se explican algunos conceptos, técnicas y tareas vinculadas a la Inteligencia Artificial, al Procesamiento de Lenguaje Natural y otros conceptos necesarios para enmarcar el proyecto. El propósito de este capítulo es contextualizar al lector en la teoría referenciada en las secciones posteriores de este informe.

En la sección ~\ref{sec:CAP2_IA} se introduce el concepto de Inteligencia Artificial. Luego, en la sección ~\ref{sec:CAP2_AA} se describen los conceptos principales asociados con el Aprendizaje Automático, donde se explican las principales etapas y técnicas utilizadas en este proyecto. En la sección ~\ref{sec:CAP2_PLN} se presenta la disciplina del Procesamiento de Lenguaje Natural, presentando los modelos en los que se basa la solución de este proyecto. Por último, en la sección ~\ref{sec:CAP2_Extra} se abordan otros conceptos útiles y complementarios para el proyecto, como las Bases de Datos y las Interfaces Web.

El lector que ya cuente con conocimientos sobre estos temas puede optar por saltear la lectura de este capítulo.

\section{Inteligencia Artificial}\label{sec:CAP2_IA}

\section{Aprendizaje Automático}\label{sec:CAP2_AA}

\subsection{Aprendizaje Supervisado}

\subsection{Aprendizaje No Supervisado}

\subsection{Entrenamiento, validación y verificación}

\subsection{Medidas de evaluación}
ACA o mas a bajo?
Porque si usamos Matrices de confucion capaz la idea seria ponerlo aca, pero si va a ser el metodo casero hasta mas abajo prodria ir porque es custom para PLN

\section{Procesamiento de Lenguaje Natural}\label{sec:CAP2_PLN}

\subsection{Preprocesamiento}

\subsubsection{Word embeddings}

\subsection{Modelos Utilizados}

\subsubsection{Latent Dirichlet Allocation}

\subsubsection{Bert Topic}

\subsubsection{Biterm Topic Modeling}

\subsubsection{Biterm Dirichlet Process}

\subsection{Medidas de evaluación}
O aca qsy

\section{Conceptos Complementarios}\label{sec:CAP2_Extra}

\subsection{Base de Datos Relacional}

\subsection{Contenedor/Docker}

\subsection{Web API}